% !TeX program = lualatex
% !TeX encoding = utf8
% !TeX spellcheck = uk_UA
% !TeX root = main.tex

%% --------------------------------------------------------
%% Додаток: шпаргалка «ключові масштаби» на одній сторінці
%% Повний список формул і пояснення -- у розд. «Основні масштаби», табл. «Основні масштаби та формули».
%% --------------------------------------------------------

\section{Шпаргалка: ключові масштаби}\label{sec:cheatsheet}\hypertarget{app-cheatsheet}{}

Практичні формули в~СГС для плазми з~концентрацією $n$ ($\text{см}^{-3}$), температурою $T$ (еВ) і~полем
$B$ (Гс); $\mu=m_i/m_p$ (для дейтерію $\mu=2$). Останній стовпець~--- значення для токамака:
$n=10^{14}~\text{см}^{-3}$, $T_e=T_i=10$~кеВ, $B=5$~Тл, дейтерій, $\ln\Lambda=17$. Докладніше про~кожну
величину~--- у~розд.~\ref{sec:plasma-scales}, табл.~\ref{tab:plasma-formulas}; перехід до~СІ~--- у~додатку~\ref{sec:SGStoSI}.

\begin{center}\small
\begin{platbl}{colspec={X[1.55,l,m]X[1.7,c,m]X[1.15,c,m]}}
	Величина & Практична формула & Токамак \\
	Плазмова частота електронів $\omega_{pe}$ & $5{,}64\times10^{4}\,\sqrt{n}$~с$^{-1}$ & $5{,}6\times10^{11}$~с$^{-1}$ \\
	Дебаївська довжина $\lambda_{De}$ & $743\,\sqrt{T_e/n}$~см & $74$~мкм \\
	Число частинок у~дебаївській сфері $N_D$ & $1{,}72\times10^{9}\,T_e^{3/2}/\sqrt{n}$ & $1{,}7\times10^{8}$ \\
	Електронна циклотронна частота $\omega_{ce}$ & $1{,}76\times10^{7}\,B$~с$^{-1}$ & $8{,}8\times10^{11}$~с$^{-1}$ \\
	Ларморівський радіус електрона $r_{Le}$ & $2{,}38\,\sqrt{T_e}/B$~см & $48$~мкм \\
	Ларморівський радіус іона $r_{Li}$ & $102\,\sqrt{\mu T_i}/B$~см & $0{,}29$~см \\
	Теплова швидкість електронів $v_{Te}$ & $4{,}19\times10^{7}\,\sqrt{T_e}$~см/с & $4{,}2\times10^{9}$~см/с \\
	Іонно-звукова швидкість $c_s$ & $9{,}79\times10^{5}\,\sqrt{T_e/\mu}$~см/с & $6{,}9\times10^{7}$~см/с \\
	Альвенівська швидкість $v_A$ & $2{,}18\times10^{11}\,B/\sqrt{\mu n}$~см/с & $7{,}7\times10^{8}$~см/с \\
	Частота зіткнень електронів $\nu_{ei}$ & $2{,}91\times10^{-6}\,n\ln\Lambda\,T_e^{-3/2}$~с$^{-1}$ & $4{,}9\times10^{3}$~с$^{-1}$ \\
	Параметр $\beta=8\pi p/B^2$ ($T_e=T_i$, $T$ в~еВ) & $8{,}05\times10^{-11}\,nT/B^2$ & $3{,}2\,\%$ \\
	Час дифузії поля $\tau_d=4\pi\sigma L^2/c^2$ ($L$ в~см) & $2{,}7\times10^{-6}\,T_e^{3/2}L^2/\ln\Lambda$~с & $1{,}6\times10^{3}$~с ($T_e=10$~кеВ, $L=1$~м) \\
\end{platbl}
\end{center}

\noindent\emphx{Чотири питання до будь-якої задачі:} який просторовий масштаб? який часовий? які процеси встигають
відбутися? яка модель (кінетика, МГД, окремі частинки) справедлива на~цих масштабах?
Порівняйте $L$ з~$\lambda_{De}$, $r_L$, $\lambda_{\text{пробігу}}$, а~$\omega$~--- з~$\omega_{p}$, $\omega_{c}$, $\nu$.
\clearpage
